% Generated draft: requires mathematical and source review.
\documentclass[11pt]{article}
\usepackage[margin=25mm]{geometry}
\usepackage{fontspec}
\setmainfont{Latin Modern Roman}
\usepackage{amsmath,amssymb,mathtools,unicode-math}
\setmathfont{Latin Modern Math}
\usepackage{enumitem,needspace,xurl,hyperref}
\setlength{\parindent}{0pt}
\setlength{\parskip}{5pt}
\setlength{\emergencystretch}{3em}
\newcommand{\N}{\mathbb{N}}
\newcommand{\Z}{\mathbb{Z}}
\newcommand{\Q}{\mathbb{Q}}
\newcommand{\R}{\mathbb{R}}
\newcommand{\C}{\mathbb{C}}
\newcommand{\F}{\mathbb{F}}
\newcommand{\A}{\mathbb{A}}
\newcommand{\PP}{\mathbb{P}}
\newcommand{\E}{\mathbb{E}}
\newcommand{\eps}{\varepsilon}
\newcommand{\dd}{\,\mathrm{d}}
\begin{document}
\section*{UnsolvedMath ID: 700464}
\textbf{Source problem code:} catalogue entry
\subsection{Definitions and mathematical statement}
Let
\[
\mathcal D_0=\{G: G\text{ is a finite group admitting a finite presentation }\langle X\mid R\rangle\text{ with }|X|=|R|\}.
\]
This is the source's definition of the class of finite groups of deficiency zero. Write \(d(G)\) for the minimum cardinality of a generating set of \(G\). For a prime \(p\), a finite \(p\)-group is a group of order \(p^n\) for some integer \(n\geq 0\).

\begin{enumerate}
\item[(a)] The source asks whether \(\mathcal D_0\) contains any \(3\)-generator \(p\)-group for \(p\geq 5\). Under the reading that \(3\)-generator means \(d(G)=3\), and that ``any'' is existential over the primes, this is the question whether
\[
\exists\,p\text{ prime},\ p\geq 5,\quad \exists\,G\in\mathcal D_0:\quad |G|=p^n\text{ for some }n\geq 0,\quad d(G)=3.
\]
The supplied wording does not explicitly distinguish minimum generator number exactly three from generation by at most three elements, nor does it specify whether existence is sought for some prime or for each prime \(p\geq 5\). These interpretations require review.

\item[(c)] For a group \(G\), define its derived series by \(G^{(0)}=G\) and \(G^{(i+1)}=[G^{(i)},G^{(i)}]\). A group is soluble if \(G^{(\ell)}=1\) for some \(\ell\geq 0\); its derived length \(\operatorname{dl}(G)\) is the least such \(\ell\). Is there an integer \(L\geq 0\) such that
\[
\forall\,G\in\mathcal D_0,\qquad G\text{ soluble}\ \Longrightarrow\ \operatorname{dl}(G)\leq L?
\]
The bound is to be uniform over all finite soluble groups in \(\mathcal D_0\).

\item[(d)] Determine, up to isomorphism, the set
\[
\left\{S:\begin{array}{l}
S\text{ is a finite non-abelian simple group, and there exist }G\in\mathcal D_0\\
\text{and a composition series }1=G_0\triangleleft G_1\triangleleft\cdots\triangleleft G_m=G\\
\text{such that }G_i/G_{i-1}\cong S\text{ for some }1\leq i\leq m
\end{array}\right\}.
\]
Here a composition series is a subnormal series whose successive quotients are nontrivial simple groups.

\item[(b), archived.] The source also records the question: are the central factors of nilpotent groups in \(\mathcal D_0\) generated by at most three elements? A group is nilpotent if its lower central series \(\gamma_1(G)=G\), \(\gamma_{i+1}(G)=[\gamma_i(G),G]\), reaches the trivial group. The expression ``central factors'' is not defined in the supplied record. If it means the central quotient \(G/Z(G)\), where
\[
Z(G)=\{z\in G: zg=gz\text{ for every }g\in G\},
\]
then the corresponding assertion is
\[
\forall\,G\in\mathcal D_0,\qquad G\text{ nilpotent}\ \Longrightarrow\ d(G/Z(G))\leq 3.
\]
This formula is conditional on that interpretation of ``central factors''; the intended quotient or family of sections requires source confirmation. The supplied archive reports a negative answer to part (b), attributed to G. Havas and E. F. Robertson (1996). No counterexample or proof is supplied here.
\end{enumerate}
The record labels the combined problem partially solved, places parts (a), (c), and (d) in the main section, and places part (b) in the solved archive. These are supplied source classifications, not independent mathematical verification.
\subsection{Short English statement}
The problem concerns finite groups admitting presentations with equally many generators and relators. It asks about three-generator \(p\)-groups for primes \(p\geq 5\), a uniform bound on soluble derived length, and possible non-abelian simple composition factors. An archived fourth question concerns generation of central factors of nilpotent groups; the supplied record reports a negative answer. The generator convention, prime quantifier in part (a), and meaning of central factors in part (b) require review.
\subsection{Sources}
\begin{itemize}
\item Supplied local UnsolvedMath registry record, permanent ID \(700464\), exact source code \texttt{catalogue entry}, title \emph{Kourovka Notebook Problem 8.12}. The record supplies the statement, background, and status-review information used in this draft.
\item \emph{The Kourovka Notebook (Unsolved Problems in Group Theory)}, 8th issue (1982), Problem 8.12, proposed by D. L. Johnson and E. F. Robertson. The supplied background identifies an update in the 21st issue (2026), \texttt{arXiv:1401.0300v48}. The collection description names E. I. Khukhro and V. D. Mazurov as editors.
\item Supplied Archive of Solved Problems attribution for part (b): G. Havas and E. F. Robertson, \emph{Commun. Algebra} \textbf{24} (1996), 3483--3487. The supplied assessment is ``No, not always.'' No article title, explicit URL, counterexample, or proof is included in the record.
\item The supplied status review is dated \texttt{2026-10-16}, labels the problem \texttt{partially\_solved}, and states that solved parts appear in the archive while remaining parts appear in the main section. This is supplied dated curation, not an independently verified literature assessment. The record identifies upstream revision \texttt{d371ee917ac9647fbb06cd6fc529deec6f3dda82}.
\end{itemize}
\end{document}
